\documentclass[11pt]{article}

\usepackage[french]{babel}

\usepackage[T1]{fontenc}

\usepackage{amsfonts}
\usepackage{amsmath}
\usepackage{amssymb}

\usepackage{graphicx}
\usepackage{multirow}
\usepackage{multicol}
\usepackage{xcolor}
\usepackage[shortlabels]{enumitem}

\usepackage[margin=2cm]{geometry}

\usepackage{mathtools}
\usepackage{siunitx}
\usepackage{physics}

\usepackage{tikz}
\usepackage{hyperref}

\usepackage{algorithm}
\usepackage{algpseudocode}
\usetikzlibrary{quantikz2,babel}

\tikzset{qcircuit/.style={
    column sep=.9cm,
    row sep=0.7cm,
    operator/.style={draw=black!80, line width=0.5pt, fill=gray!5, minimum height=0.7cm, minimum width=0.7cm},
    phase/.style={draw=black!80, fill=black!20, circle, minimum size=4pt},
    slice/.style={font=\footnotesize\itshape, color=gray!50, dashed}
}}


\DeclarePairedDelimiter\ceil{\lceil}{\rceil}
\DeclarePairedDelimiter\floor{\lfloor}{\rfloor}


\AtBeginDocument{\RenewCommandCopy\qty\SI}

\makeatletter
   \def\vhrulefill#1{\leavevmode\leaders\hrule\@height#1\hfill \kern\z@}
\makeatother

\newcommand{\moy}[1]{\langle #1 \rangle}
\newcommand{\pDer}[2]{\frac{\partial #1}{\partial #2}}
\newcommand{\pDerd}[2]{\frac{\partial^2 #1}{\partial #2^2}}

\newcounter{questioncounter}
\newenvironment{question}[1]{
\refstepcounter{questioncounter}
\noindent\rule{0.05\textwidth}{1pt}\,\,\raisebox{-.3\ht\strutbox}{\textbf{Question \thequestioncounter \hspace{5mm} #1}}\,\,\vhrulefill{1pt}\par\vspace{0.4cm}
}{\par}

%======================================================
%en-tete
%======================================================
\begin{document}
\selectlanguage{french}
\noindent{\Large Anthony Wroblewski} \hfill {\Large \today}\\
%{\Large Anthony Wroblewski}

\par\vspace{5mm}
\begin{center}
   {\Large IFT436 - Algorithmes et structures de données} \\ \vspace{5mm}
   {\LARGE\sffamily Devoir \#2} \\ \vspace{5mm}
\end{center}

%======================================================
%question 1
%======================================================
\begin{question}{Une question d’échauffement}

    
\end{question}


%======================================================
%question 2
%======================================================
\begin{question}{Le problème de la sous-somme}

\begin{enumerate}[label=(\alph*)]

\item 
\end{enumerate}
\end{question}
%======================================================
%question 3
%======================================================
\begin{question}{Puissances en temps logarithmique}
    \begin{enumerate}[label=(\alph*)]
\item  
\end{enumerate}
\end{question}


\bibliographystyle{apsrev4-2}
\bibliography{references}
\end{document}
